% Options for packages loaded elsewhere
\PassOptionsToPackage{unicode}{hyperref}
\PassOptionsToPackage{hyphens}{url}
\documentclass[
  ignorenonframetext,
]{beamer}
\newif\ifbibliography
\usepackage{pgfpages}
\setbeamertemplate{caption}[numbered]
\setbeamertemplate{caption label separator}{: }
\setbeamercolor{caption name}{fg=normal text.fg}
\beamertemplatenavigationsymbolsempty
% remove section numbering
\setbeamertemplate{part page}{
  \centering
  \begin{beamercolorbox}[sep=16pt,center]{part title}
    \usebeamerfont{part title}\insertpart\par
  \end{beamercolorbox}
}
\setbeamertemplate{section page}{
  \centering
  \begin{beamercolorbox}[sep=12pt,center]{section title}
    \usebeamerfont{section title}\insertsection\par
  \end{beamercolorbox}
}
\setbeamertemplate{subsection page}{
  \centering
  \begin{beamercolorbox}[sep=8pt,center]{subsection title}
    \usebeamerfont{subsection title}\insertsubsection\par
  \end{beamercolorbox}
}
% Prevent slide breaks in the middle of a paragraph
\widowpenalties 1 10000
\raggedbottom
\AtBeginPart{
  \frame{\partpage}
}
\AtBeginSection{
  \ifbibliography
  \else
    \frame{\sectionpage}
  \fi
}
\AtBeginSubsection{
  \frame{\subsectionpage}
}
\usepackage{iftex}
\ifPDFTeX
  \usepackage[T1]{fontenc}
  \usepackage[utf8]{inputenc}
  \usepackage{textcomp} % provide euro and other symbols
\else % if luatex or xetex
  \usepackage{unicode-math} % this also loads fontspec
  \defaultfontfeatures{Scale=MatchLowercase}
  \defaultfontfeatures[\rmfamily]{Ligatures=TeX,Scale=1}
\fi
\usepackage{lmodern}
\ifPDFTeX\else
  % xetex/luatex font selection
\fi
% Use upquote if available, for straight quotes in verbatim environments
\IfFileExists{upquote.sty}{\usepackage{upquote}}{}
\IfFileExists{microtype.sty}{% use microtype if available
  \usepackage[]{microtype}
  \UseMicrotypeSet[protrusion]{basicmath} % disable protrusion for tt fonts
}{}
\makeatletter
\@ifundefined{KOMAClassName}{% if non-KOMA class
  \IfFileExists{parskip.sty}{%
    \usepackage{parskip}
  }{% else
    \setlength{\parindent}{0pt}
    \setlength{\parskip}{6pt plus 2pt minus 1pt}}
}{% if KOMA class
  \KOMAoptions{parskip=half}}
\makeatother
\usepackage{color}
\usepackage{fancyvrb}
\newcommand{\VerbBar}{|}
\newcommand{\VERB}{\Verb[commandchars=\\\{\}]}
\DefineVerbatimEnvironment{Highlighting}{Verbatim}{commandchars=\\\{\}}
% Add ',fontsize=\small' for more characters per line
\newenvironment{Shaded}{}{}
\newcommand{\AlertTok}[1]{\textcolor[rgb]{1.00,0.00,0.00}{\textbf{#1}}}
\newcommand{\AnnotationTok}[1]{\textcolor[rgb]{0.38,0.63,0.69}{\textbf{\textit{#1}}}}
\newcommand{\AttributeTok}[1]{\textcolor[rgb]{0.49,0.56,0.16}{#1}}
\newcommand{\BaseNTok}[1]{\textcolor[rgb]{0.25,0.63,0.44}{#1}}
\newcommand{\BuiltInTok}[1]{\textcolor[rgb]{0.00,0.50,0.00}{#1}}
\newcommand{\CharTok}[1]{\textcolor[rgb]{0.25,0.44,0.63}{#1}}
\newcommand{\CommentTok}[1]{\textcolor[rgb]{0.38,0.63,0.69}{\textit{#1}}}
\newcommand{\CommentVarTok}[1]{\textcolor[rgb]{0.38,0.63,0.69}{\textbf{\textit{#1}}}}
\newcommand{\ConstantTok}[1]{\textcolor[rgb]{0.53,0.00,0.00}{#1}}
\newcommand{\ControlFlowTok}[1]{\textcolor[rgb]{0.00,0.44,0.13}{\textbf{#1}}}
\newcommand{\DataTypeTok}[1]{\textcolor[rgb]{0.56,0.13,0.00}{#1}}
\newcommand{\DecValTok}[1]{\textcolor[rgb]{0.25,0.63,0.44}{#1}}
\newcommand{\DocumentationTok}[1]{\textcolor[rgb]{0.73,0.13,0.13}{\textit{#1}}}
\newcommand{\ErrorTok}[1]{\textcolor[rgb]{1.00,0.00,0.00}{\textbf{#1}}}
\newcommand{\ExtensionTok}[1]{#1}
\newcommand{\FloatTok}[1]{\textcolor[rgb]{0.25,0.63,0.44}{#1}}
\newcommand{\FunctionTok}[1]{\textcolor[rgb]{0.02,0.16,0.49}{#1}}
\newcommand{\ImportTok}[1]{\textcolor[rgb]{0.00,0.50,0.00}{\textbf{#1}}}
\newcommand{\InformationTok}[1]{\textcolor[rgb]{0.38,0.63,0.69}{\textbf{\textit{#1}}}}
\newcommand{\KeywordTok}[1]{\textcolor[rgb]{0.00,0.44,0.13}{\textbf{#1}}}
\newcommand{\NormalTok}[1]{#1}
\newcommand{\OperatorTok}[1]{\textcolor[rgb]{0.40,0.40,0.40}{#1}}
\newcommand{\OtherTok}[1]{\textcolor[rgb]{0.00,0.44,0.13}{#1}}
\newcommand{\PreprocessorTok}[1]{\textcolor[rgb]{0.74,0.48,0.00}{#1}}
\newcommand{\RegionMarkerTok}[1]{#1}
\newcommand{\SpecialCharTok}[1]{\textcolor[rgb]{0.25,0.44,0.63}{#1}}
\newcommand{\SpecialStringTok}[1]{\textcolor[rgb]{0.73,0.40,0.53}{#1}}
\newcommand{\StringTok}[1]{\textcolor[rgb]{0.25,0.44,0.63}{#1}}
\newcommand{\VariableTok}[1]{\textcolor[rgb]{0.10,0.09,0.49}{#1}}
\newcommand{\VerbatimStringTok}[1]{\textcolor[rgb]{0.25,0.44,0.63}{#1}}
\newcommand{\WarningTok}[1]{\textcolor[rgb]{0.38,0.63,0.69}{\textbf{\textit{#1}}}}
\setlength{\emergencystretch}{3em} % prevent overfull lines
\providecommand{\tightlist}{%
  \setlength{\itemsep}{0pt}\setlength{\parskip}{0pt}}
% Slide layout defaults for warning-clean scientific decks.
\usepackage{etoolbox}
\IfFileExists{xurl.sty}{\usepackage{xurl}}{}
\IfFileExists{seqsplit.sty}{\usepackage{seqsplit}}{\newcommand{\seqsplit}[1]{#1}}
\protected\def\breakseq#1{\seqsplit{#1}}
\protected\def\breaktt#1{\begingroup\ttfamily\seqsplit{#1}\endgroup}
\setlength{\emergencystretch}{6em}
\tolerance=5000
\hbadness=10000
\hfuzz=1pt
\setlength{\tabcolsep}{2pt}
\AtBeginEnvironment{longtable}{\tiny\renewcommand{\arraystretch}{0.86}\setlength{\tabcolsep}{1pt}}
\AtBeginEnvironment{tabular}{\tiny\renewcommand{\arraystretch}{0.86}\setlength{\tabcolsep}{1pt}}

% Natbib and cross-reference fallbacks — slides don't load natbib
% or cleveref, but combined-PDF manuscript prose may emit these
% commands. The fallback renders citations as a bracketed key list
% and cross-references as detokenized labels so slides don't fail on
% undefined control sequences or raw underscores. \providecommand is
% a no-op if packages load later.
\providecommand{\citep}[1]{[#1]}
\providecommand{\citet}[1]{#1}
\providecommand{\citealp}[1]{#1}
\providecommand{\citeauthor}[1]{#1}
\providecommand{\citeyear}[1]{#1}
\providecommand{\cref}[1]{\texttt{\detokenize{#1}}}
\providecommand{\Cref}[1]{\texttt{\detokenize{#1}}}
\usepackage{bookmark}
\IfFileExists{xurl.sty}{\usepackage{xurl}}{} % add URL line breaks if available
\urlstyle{same}
% fallback for those not using the hyperref driver hyperxmp:
\makeatletter
\@ifundefined{xmpquote}{\newcommand{\xmpquote}[1]{#1}}{}
\makeatother
\hypersetup{
  hidelinks,
  pdfcreator={LaTeX via pandoc}}

\author{\texorpdfstring{}{}}
\date{}

\begin{document}

\section{Appendix A: Tooling and
Reproduction}\label{appendix-a-tooling-and-reproduction}

\begin{frame}[fragile,allowframebreaks]{Appendix A: Tooling and
Reproduction}
The pipeline is a two-layer system: generic infrastructure (rendering,
validation, logging) shared across the template monorepo, and
project-local \texttt{src/} modules that implement the meta-analysis.
All numbered \texttt{scripts/} are thin orchestrators that wire I/O,
configuration loading, and logging --- no computational logic resides in
scripts.
\end{frame}

\begin{frame}[fragile,allowframebreaks]{Reproduce the Offline Default
Run}
\protect\phantomsection\label{reproduce-the-offline-default-run}
No network, no language model required:

\begin{Shaded}
\begin{Highlighting}[]
\ExtensionTok{uv}\NormalTok{ run python scripts/generate\_fixture\_corpus.py }\AttributeTok{{-}{-}out}\NormalTok{ output/data/corpus.jsonl}
\ExtensionTok{uv}\NormalTok{ run python scripts/02\_meta\_analysis\_pipeline.py}
\ExtensionTok{uv}\NormalTok{ run python scripts/03\_build\_knowledge\_graph.py }\AttributeTok{{-}{-}max{-}papers}\NormalTok{ 0}
\ExtensionTok{uv}\NormalTok{ run python scripts/04\_generate\_figures.py }\AttributeTok{{-}{-}dpi}\NormalTok{ 300}
\ExtensionTok{uv}\NormalTok{ run python scripts/05\_inject\_variables.py}
\end{Highlighting}
\end{Shaded}
\end{frame}

\begin{frame}[fragile,allowframebreaks]{Reproduce the Live Run}
\protect\phantomsection\label{reproduce-the-live-run}
This manuscript was generated from a live retrieval run. To reproduce:

\begin{Shaded}
\begin{Highlighting}[]
\CommentTok{\# Live search (all 7 engines, max 1000 per engine)}
\ExtensionTok{uv}\NormalTok{ run python scripts/01\_literature\_search.py }\AttributeTok{{-}{-}query}\NormalTok{ modafinil }\AttributeTok{{-}{-}max{-}results}\NormalTok{ 1000 }\AttributeTok{{-}{-}no{-}resume}

\CommentTok{\# Analysis pipeline}
\ExtensionTok{uv}\NormalTok{ run python scripts/02\_meta\_analysis\_pipeline.py}
\ExtensionTok{uv}\NormalTok{ run python scripts/03\_build\_knowledge\_graph.py }\AttributeTok{{-}{-}max{-}papers}\NormalTok{ 0}
\ExtensionTok{uv}\NormalTok{ run python scripts/04\_generate\_figures.py }\AttributeTok{{-}{-}dpi}\NormalTok{ 300}
\ExtensionTok{uv}\NormalTok{ run python scripts/05\_inject\_variables.py}
\ExtensionTok{uv}\NormalTok{ run python scripts/06\_fulltext\_assessment.py}
\end{Highlighting}
\end{Shaded}
\end{frame}

\begin{frame}[fragile,allowframebreaks]{Re-target to Another Topic}
\protect\phantomsection\label{re-target-to-another-topic}
Edit \breaktt{manuscript/config.yaml} ---
\breaktt{project\_config.search.term}, \texttt{query},
\breaktt{relevance\_keywords}, \breaktt{subfield\_keywords}, and
\breaktt{hypothesis\_definitions} --- then regenerate the seed corpus and
re-run. No code changes are required; the manuscript re-targets through
token injection.
\end{frame}

\begin{frame}[fragile,allowframebreaks]{Live Retrieval}
\protect\phantomsection\label{live-retrieval}
Enable engines under \breaktt{project\_config.search.engines}, supply any
optional credentials (Unpaywall email, Semantic Scholar key), and run
\breaktt{scripts/01\_literature\_search.py}; absent engines degrade to
skipped sources. The CLI supports per-engine skip flags:
\texttt{-\/-skip-arxiv}, \texttt{-\/-skip-s2},
\breaktt{-\/-skip-openalex}, \breaktt{-\/-skip-crossref},
\texttt{-\/-skip-pubmed}, \breaktt{-\/-skip-sovietrxiv},
\breaktt{-\/-skip-chinarxiv}.
\end{frame}

\begin{frame}[fragile,allowframebreaks]{Test Suite}
\protect\phantomsection\label{test-suite}
Every stage is covered by a no-mocks test suite (real computation and
\breaktt{pytest-httpserver} for network adapters) gated at \(\geq 90\%\)
statement coverage on \texttt{src/}. The suite includes 819 tests
covering:

\begin{itemize}
\tightlist
\item
  Record models and serialization (deduplication, canonical ID
  hierarchy)
\item
  All 7 engine clients (arXiv, Semantic Scholar, OpenAlex, Crossref,
  PubMed, SovietRxiv, ChinaRxiv) with pytest-httpserver integration
  tests
\item
  Search runner (multi-engine dispatch, relevance filtering,
  resume/clear, YAML config)
\item
  Bibliometric analysis (subfield classification, temporal metrics,
  TF-IDF, NMF, citation network)
\item
  Knowledge graph (schema, nanopublications, hypothesis scoring, LLM
  extraction)
\item
  Visualization (headless figure generation, style config)
\item
  Manuscript variable computation and injection
\end{itemize}
\end{frame}

\end{document}
